\documentclass[12pt, a4paper]{article}

% Paquetes requeridos para el circuito
\usepackage[utf8]{inputenc}
\usepackage[T1]{fontenc}
\usepackage{circuitikz} % Paquete principal para dibujar circuitos

\begin{document}

\begin{center}
    % El entorno circuitikz genera el gráfico. 
    % Usamos 'american' para los símbolos estándar de resistencias (zig-zag).
    \begin{circuitikz}[american]
        
        % 1. Fuente V1 (de (0,0) a (0,2))
        \draw (0,0) to[V, l=$V_1$] (0,2);
        
        % 2. Resistencia R1 en paralelo (de (2,0) a (2,2))
        \draw (2,0) to[R, l=$R_1$] (2,2);
        
        % Conexiones superior e inferior para completar el paralelo
        \draw (0,2) -- (2,2);
        \draw (0,0) -- (2,0);
        
        % 3. Resistencia R2 en serie saliendo de la unión superior hacia el nodo A
        \draw (2,2) to[R, l=$R_2$, -o] (4.5,2) node[right] {A};
        
        % 4. Terminal B saliendo de la unión inferior
        \draw (2,0) to[short, -o] (4.5,0) node[right] {B};
        
    \end{circuitikz}
\end{center}

\end{document}